\subsection{Applications : II. The multiplicative group of units of the integers modulo a}

Let \(a\) be a rational integer \(> 1\) , and let \((\mathbf{Z} / a\mathbf{Z})^*\) be the multiplicative group of invertible elements of the ring \(\mathbf{Z} / a\mathbf{Z}\) . If \(a = \prod_{i} p_i^{n(i)}\) is the decomposition of \(a\) into

prime factors, then the ring \(\mathbf{Z}/a\mathbf{Z}\) is isomorphic to the product of the rings \(\mathbf{Z}/p_{i}^{n(i)}\mathbf{Z}\) (VII, p.~3, Prop. 4) and the group \((\mathbf{Z}/a\mathbf{Z})^{*}\) is isomorphic to the product of the groups \((\mathbf{Z}/p_{i}^{n(i)}\mathbf{Z})^{*}\) . We are thus reduced to the study of the groups \((\mathbf{Z}/p^{n}\mathbf{Z})^{*}\) , where \(p\) is a prime number; recall (V, p.~80) that the order \(\varphi(p^{n})\) of \((\mathbf{Z}/p^{n}\mathbf{Z})^{*}\) is \(p^{n}-p^{n-1}=p^{n-1}(p-1)\) .

Suppose first of all that \(p > 2\) ; the natural homomorphism \(\mathbf{Z} / p^n\mathbf{Z} \to \mathbf{Z} / p\mathbf{Z}\) restricts to a homomorphism of groups from \((\mathbf{Z} / p^n\mathbf{Z})^*\) onto \((\mathbf{Z} / p\mathbf{Z})^*\) , whose kernel we denote \(\mathrm{U}(p^n)\) ; by VII, p.~3, Prop. 3 the residue class mod \(p^n\) of an integer \(m\) is invertible if and only if \(m\) is coprime to \(p\) , that is if and only if the residue class of \(m \mod p\) is invertible. It follows that \(\mathrm{U}(p^n)\) consists of all the residue classes mod \(p^n\) of integers congruent to \(1 \mod p\) , so has \(p^{n-1}\) elements, and that there is an exact sequence

\[
\{1 \} \rightarrow \mathrm{U} (p ^ {n}) \rightarrow (\mathbf {Z} / p ^ {n} \mathbf {Z}) ^ {*} \rightarrow (\mathbf {Z} / p \mathbf {Z}) ^ {*} \rightarrow \{1 \}.\tag{3}
\]

Similarly, for \(n \geqslant 2\) let \(U(2^n)\) denote the kernel of the natural homomorphism from \((\mathbf{Z}/2^n\mathbf{Z})^*\) to \((\mathbf{Z}/4\mathbf{Z})^*\) ; this is a group of order \(2^{n-2}\) , consisting of all the residue classes mod \(2^n\) of integers congruent to 1 mod 4, and there is an exact sequence

\[
\{1 \} \rightarrow \mathrm{U} (2 ^ {n}) \rightarrow (\mathbf {Z} / 2 ^ {n} \mathbf {Z}) ^ {*} \rightarrow (\mathbf {Z} / 4 \mathbf {Z}) ^ {*} \rightarrow \{1 \}.\tag{4}
\]

Lemma 3. --- Let \(x, y, k\) be integers with \(k \geqslant 0\) , and let \(p > 2\) be a prime number. If \(x \equiv 1 + py \mod p^2\) then \(x^{p^k} \equiv 1 + p^{k+1}y \mod p^{k+2}\) . If \(x \equiv 1 + 4y \mod 8\) then \(x^{2^k} \equiv 1 + 2^{k+2}y \mod 2^{k+3}\) .

To prove the first assertion, it is enough to show that, if \(k \geqslant 1\) and \(x \equiv 1 + p^{k}y \mod p^{k+1}\) , then \(x^{p} \equiv 1 + p^{k+1}y \mod p^{k+2}\) , and then to argue by induction on the integer k. For all \(a \in Z\) and \(k \geqslant 1\) , it is immediate that

\[
(1 + p ^ {k} a) ^ {p} \equiv 1 + p ^ {k + 1} a \mod p ^ {k + 2},
\]

hence

\[
\begin{array}{r l} (1 + p ^ {k} y + p ^ {k + 1} z) ^ {p} & = (1 + p ^ {k} (y + p z)) ^ {p} \equiv \\ & \equiv 1 + p ^ {k + 1} (y + p z) \equiv 1 + p ^ {k + 1} y \mod p ^ {k + 2} \end{array}
\]

Similarly, for \(k \geqslant 1\) we have

\[
(1 + 2 ^ {k + 1} a) ^ {2} \equiv 1 + 2 ^ {k + 2} a \mod 2 ^ {k + 3},
\]

so

\[
(1 + 2 ^ {k + 1} y + 2 ^ {k + 2} z) ^ {2} \equiv 1 + 2 ^ {k + 2} y \mod 2 ^ {k + 3},
\]

whence the second assertion by induction on k.

PROPOSITION 3. --- Let p \textgreater{} 2 be a prime number and let n \textgreater{} 0 be an integer ; then the group U(p\^{}n) is cyclic of order p\^{}\{n-1\} ; if n ≥ 2 then the residue class mod p\^{}n of an integer x congruent to 1 mod p is a generator of U(p\^{}n) if and only if x is not congruent to 1 mod p\^{}2. Let m \textgreater{} 1 be an integer ; then the group U(2\^{}m) is cyclic of order 2\^{}\{m-2\} ; if m ≥ 3 then the residue class mod 2\^{}m of an integer x congruent to 1 mod 4 is a generator of U(2\^{}m) if and only if x is not congruent to 1 mod 8.

Since \(\mathrm{U}(p^{n})\) has order \(p^{n-1}\) , the order of every element u of \(\mathrm{U}(p^{n})\) is a power of p, and u is a generator of \(\mathrm{U}(p^{n})\) if and only if \(u^{p^{n-2}} \neq 1\) . Now if u is the class of \(x = 1 + py\) , then \(u^{p^{n-2}}\) is the class of \(1 + p^{n-1}y\) , by Lemma 3, whence u generates \(\mathrm{U}(p^{n})\) if and only if \(y \neq 0 \mod p\) , in other words \(x \neq 1 \mod p^{2}\) . For example, the class \(1 + p\) generates \(\mathrm{U}(p^{n})\) . Similarly, the class u of x mod \(2^{n}\) generates \(\mathrm{U}(2^{n})\) if and only if \(u^{2^{n-3}} \neq 1\) , which means that x is not congruent to 1 mod 8, by Lemma 3; this is satisfied by x = 5.

Lemma 4. --- Let A be a principal ideal domain and let \(0 \rightarrow N \rightarrow M \rightarrow P \rightarrow 0\) be an exact sequence of A-modules. Suppose that there exist coprime elements a, \(b \in A\) such that aN = 0 and bP = 0. Then the exact sequence splits. If in addition N and P are both cyclic, then M is cyclic.

The module M is torsion, since abM = 0. The first assertion follows from Cor. 3 of VII, p.~9. If N and P are cyclic, then they are finitely generated, and hence so is M (II, p.~17, Cor. 5); since each p-primary component of M is isomorphic to a p-primary component either of N or of P, it follows from Remark 2 of VII, p.~10, that M is cyclic.

THEOREM 3. --- If \(a = \prod_{i} p_{i}^{n(i)}\) is the prime decomposition of the integer

a \textgreater{} 1, then the group \((\mathbf{Z}/a\mathbf{Z})^{*}\) of invertible elements of the ring Z/aZ is isomorphic to the product of the groups \((\mathbf{Z}/p_{i}^{n(i)}\mathbf{Z})^{*}\) . If p \textgreater{} 2 is a prime number and \(n \geqslant 1\) an integer, then the group \((\mathbf{Z}/p^{n}\mathbf{Z})^{*}\) is cyclic of order \(p^{n-1}(p - 1)\) . The group \((\mathbf{Z}/2\mathbf{Z})^{*}\) is trivial; for \(n \geqslant 2\) the group \((\mathbf{Z}/2^{n}\mathbf{Z})^{*}\) is the direct product of the cyclic group of order \(2^{n-2}\) generated by the residue class of 5 mod \(2^{n}\) and the cyclic group of order 2 consisting of the residue classes of 1 and -1 mod \(2^{n}\) .

The orders \(p^{n-1}\) of \(U(p^n)\) and \(p-1\) of \((\mathbf{Z}/p\mathbf{Z})^*\) are coprime; since \(U(p^n)\) and \((\mathbf{Z}/p\mathbf{Z})^*\) are cyclic (Prop. 3 and V, p.~78, Lemma 1), the group \((\mathbf{Z}/p^n\mathbf{Z})^*\) is cyclic (apply Lemma 4 to the exact sequence (3)). If \(n \geqslant 2\) then the restriction of the homomorphism \(v: (\mathbf{Z}/2^n\mathbf{Z})^* \to (\mathbf{Z}/4\mathbf{Z})^*\) to the subgroup \(\{1, -1\}\) is bijective; the group \((\mathbf{Z}/2^n\mathbf{Z})^*\) is thus the direct product of this subgroup and the kernel \(U(2^n)\) of \(v\) ; the result follows from Prop. 3.

Remark. --- Let p \textgreater{} 2 be a prime number and let x be an integer congruent to 1 mod p and not congruent to 1 mod \(p^{2}\) ; there is an exact sequence

\[
\{0 \} \rightarrow \mathbf {Z} / p ^ {n - 1} \mathbf {Z} \xrightarrow {u} (\mathbf {Z} / p ^ {n} \mathbf {Z}) ^ {*} \xrightarrow {v} (\mathbf {Z} / p \mathbf {Z}) ^ {*} \rightarrow \{1 \}\tag{5}
\]

of groups, where v is the natural projection and where u is induced on the quotients by the map \(r \mapsto x^{r}\) . Let \(Z_{p}\) be the ring of p-adic integers (V, p.~96), and let x be an element of \(Z_{p}\) such that \(x - 1 \in pZ_{p}\) and \(x - 1 \notin p^{2}Z_{p}\) ; on passing to inverse limits, the exact sequences (5) induce an exact sequence

\[
\{0 \} \rightarrow \mathbf {Z} _ {p} \stackrel {u} {\rightarrow} \mathbf {Z} _ {p} ^ {*} \stackrel {v} {\rightarrow} (\mathbf {Z} / p \mathbf {Z}) ^ {*} \rightarrow \{1 \}
\]

where \(v\) is the natural map, and the continuous map \(u\) extends the map \(n \mapsto x^n\) ( \(n \in \mathbf{Z}\) ). We often put \(x^n = u(x)\) for \(n \in \mathbf{Z}_p\) .

Similarly, if \(x \in \mathbf{Z}_2\) with \(x - 1 \in 4\mathbf{Z}_2\) and \(x - 1 \notin 8\mathbf{Z}_2\) , then there is a split exact sequence

\[
\{0 \} \rightarrow \mathbf {Z} _ {2} \stackrel {{u}} {{\rightarrow}} \mathbf {Z} _ {2} ^ {*} \stackrel {{v}} {{\rightarrow}} (\mathbf {Z} / 4 \mathbf {Z}) ^ {*} \rightarrow \{1 \},
\]

where \(u\) is a continuous extension of the map \(n \mapsto x^n\) .

